% Source: physical PDF pp. 459--468 (printed pp. 436--445). Coverage:
% complete Section 23.3, from its heading through immediately before
% Section 23.4; Eqs. (23.3.1)--(23.3.25), five unnumbered displays,
% one centered divider, citations [2], [10], and [12]--[21], and no
% figures, tables, or footnotes.
\subsection{Monopoles}
\label{sec:23.3}

As a detailed example of a topologically non-trivial field configuration,
we shall now consider the monopole of 't~Hooft and
Polyakov~\hyperref[ch23-ref-2]{[2]}, and its generalizations. We saw in
\hyperref[sec:23.1]{Section 23.1} that when a simply connected gauge group
$G$ is spontaneously broken to the $U(1)$ of electromagnetism, the
configurations of finite energy are classified according to the elements
of the group
$\pi_2(G/U(1))=\pi_1(U(1))=\mathbb Z$. (The case of non-simply connected
Lie groups is considered at the end of this section.) According to the
physical interpretation of homotopy groups discussed in
\hyperref[sec:23.2]{Section 23.2}, this means that these configurations
have a conserved additive quantum number. But we still need to show that
any of these stationary configurations actually exist, and to give a
physical interpretation of their topological quantum numbers.

As an illustrative example, consider a theory (like the Georgi--Glashow
electroweak model~\hyperref[ch23-ref-12]{[12]}) in which an $SU(2)$ gauge
group is spontaneously broken by the vacuum expectation value of an
$SU(2)$ triplet of scalar fields $\phi_n$. (It is explained at the end of
this section why in this case we say that the gauge group is $SU(2)$
rather than $SO(3)$.) The Lagrangian density for the scalars and gauge
fields in Minkowskian spacetime is taken as
\begin{equation}
 \mathcal L=-\frac14F_{n\mu\nu}F_n^{\mu\nu}
 -\frac12D_\mu\phi_nD^\mu\phi_n-V(\phi_n\phi_n),
 \tag{23.3.1}\label{eq:23.3.1}
\end{equation}
where
\begin{equation}
 F_{n\mu\nu}\equiv
 \partial_\mu A_{n\nu}-\partial_\nu A_{n\mu}
 +e\epsilon_{nml}A_{m\mu}A_{l\nu},
 \tag{23.3.2}\label{eq:23.3.2}
\end{equation}
\begin{equation}
 D_\mu\phi_n\equiv
 \partial_\mu\phi_n+e\epsilon_{nml}A_{m\mu}\phi_l,
 \tag{23.3.3}\label{eq:23.3.3}
\end{equation}
and the function $V(\phi_n\phi_n)$ is assumed to be positive, with the
value zero at a non-vanishing value $\langle\phi\rangle$ (taken positive)
of $\sqrt{\phi_n\phi_n}$. (In much work on monopoles $V$ is taken to be
the quartic polynomial
$\lambda(\phi_n\phi_n-\langle\phi\rangle^2)^2$, with $\lambda>0$, but we
shall not make this assumption here.) Eq.~\eqref{eq:23.3.1} describes a
theory with a spacetime-independent vacuum solution with $A_{n\mu}=0$,
in which a vacuum expectation value $\phi_n$ with
$\phi_n\phi_n=\langle\phi\rangle^2$ breaks $SU(2)$ to its $U(1)$
subgroup, which can be identified with the gauge group of electrodynamics.
Instead we here shall seek topologically non-trivial inhomogeneous but
time-independent classical solutions in temporal gauge, with $A_n^0=0$
but $A_n^i\ne0$. The Lagrangian density in this case is the negative of
the potential energy density $\mathcal H$, given by
\begin{equation}
 \mathcal H=\frac14F_{nij}^{\,2}
 +\frac12(D_i\phi_n)^2+V(\phi),
 \tag{23.3.4}\label{eq:23.3.4}
\end{equation}
with squares including obvious index contractions. Because each term in
\eqref{eq:23.3.4} is positive, the integral of each term must separately
converge in a configuration of finite energy.

In particular, for the integral of $V(\phi_n\phi_n)$ to converge, the
vector $\phi_n$ must have the fixed length $\langle\phi\rangle$ at
infinity, so each configuration of finite energy defines a smooth mapping
of a large two-sphere $\mathcal S$ surrounding the monopole configuration
into the two-sphere of the $\phi_n$ with
$\phi_n\phi_n=\langle\phi\rangle^2$. As $\hat{\mathbf x}$ runs over
$\mathcal S$, $\phi_n$ may run over the sphere
$\phi_n\phi_n=\langle\phi\rangle^2$ any integer number $N$ of times,
either with Jacobian
$\operatorname{Det}(\partial\mathbf x/\partial\boldsymbol\phi)$ positive,
in which case we say that the winding number is $N$, or with Jacobian
negative, in which case the winding number is $-N$.

To see what the winding number has to do with the magnetic monopole
moment, we must first consider what in this theory is observed as `the'
magnetic field. Whatever the field configuration, we can introduce a
gauge in which the scalar field $\phi_n$ points in some definite direction,
say the three-direction, in any given finite region, so that in this region
the gauge field associated with the unbroken $U(1)$ subgroup of $SU(2)$ is
$A_{3i}$. 't~Hooft~\hyperref[ch23-ref-2]{[2]} found a gauge-invariant
tensor which reduces to the usual electromagnetic field strength tensor
$\partial_\mu A_{3\nu}-\partial_\nu A_{3\mu}$ in this gauge:
\begin{equation}
 \mathcal F_{\mu\nu}\equiv
 F_{n\mu\nu}\hat\phi_n
 -\frac1e\epsilon_{nml}\hat\phi_n
 D_\mu\hat\phi_mD_\nu\hat\phi_l,
 \tag{23.3.5}\label{eq:23.3.5}
\end{equation}
where
$\hat\phi_n\equiv\phi_n/\sqrt{\phi_m\phi_m}$. To check that
$\mathcal F_{\mu\nu}$ is the ordinary electromagnetic field strength
tensor in a gauge with constant $\hat\phi_n$ (and for later purposes) we
use Eqs.~\eqref{eq:23.3.2} and \eqref{eq:23.3.3} and the identity
$\epsilon_{abc}\epsilon_{ade}
=\delta_{bd}\delta_{ce}-\delta_{be}\delta_{cd}$ to write
$\mathcal F_{\mu\nu}$ in the form~\hyperref[ch23-ref-13]{[13]}
\begin{equation}
 \mathcal F_{\mu\nu}
 =\partial_\mu(\hat\phi_nA_{n\nu})
 -\partial_\nu(\hat\phi_nA_{n\mu})
 -\frac1e\epsilon_{nml}\hat\phi_n
 \partial_\mu\hat\phi_m\partial_\nu\hat\phi_l.
 \tag{23.3.6}\label{eq:23.3.6}
\end{equation}
Thus in a gauge in which $\hat\phi_n$ is a fixed unit vector in the
three-direction, we have as promised
\[
 \mathcal F_{\mu\nu}
 =\partial_\mu A_{3\nu}-\partial_\nu A_{3\mu}.
\]

The magnetic monopole moment $g$ of any localized field configuration is
defined as $1/4\pi$ times the magnetic flux through a large closed surface
$\mathcal S$ around the configuration:
\begin{equation}
 4\pi g\equiv\frac12\epsilon_{ijk}
 \int_{\mathcal S}\mathcal F_{ij}\,d^2S_k.
 \tag{23.3.7}\label{eq:23.3.7}
\end{equation}
The first two terms in Eq.~\eqref{eq:23.3.6} for $\mathcal F_{ij}$ are
derivatives, and therefore do not contribute to the integral
\eqref{eq:23.3.7}, so that
\begin{equation}
 g=-\frac1{8\pi e}\epsilon_{ijk}\epsilon_{nml}
 \int_{\mathcal S}\hat\phi_n\partial_i\hat\phi_m
 \partial_j\hat\phi_l\,d^2S_k.
 \tag{23.3.8}\label{eq:23.3.8}
\end{equation}
This has the important property of being a topological invariant:
integrating by parts where necessary, we can see that the change in $g$
under an infinitesimal variation $\delta\hat\phi_n$ in $\hat\phi_n$ is
\[
 \delta g=-\frac3{8\pi e}\epsilon_{ijk}\epsilon_{nml}
 \int_{\mathcal S}\delta\hat\phi_n\partial_i\hat\phi_m
 \partial_j\hat\phi_l\,d^2S_k.
\]
But because $\hat\phi$ is a unit vector, $\delta\hat\phi$ as well as
$\partial_i\hat\phi$ and $\partial_j\hat\phi$ are all in the plane
perpendicular to $\hat\phi$, so
\[
 \epsilon_{nml}\delta\hat\phi_n
 \partial_i\hat\phi_m\partial_j\hat\phi_l=0,
\]
and therefore $\delta g=0$. The quantity \eqref{eq:23.3.8} is related to
the topological invariant known as the \emph{Kronecker index}.

Because $g$ is a surface integral, it is additive: for any two distant
localized configurations, the surface $\mathcal S$ used to calculate $g$
may be taken to be a pair of spheres, one surrounding each configuration,
connected by a thin neck between them, so the value of $g$ for the whole
system will be the sum of the values of $g_1$, $g_2$ for the individual
localized configurations. Furthermore, since $g$ is a topological
invariant, we will have $g=g_1+g_2$ for any field configuration that is
formed by a smooth fusion of two configurations with magnetic monopole
moments $g_1$ and $g_2$. It follows that $g$ must be proportional to the
winding number. Arafune, Freund, and Goebel~\hyperref[ch23-ref-13]{[13]}
verified this and calculated the coefficient of proportionality by using
the formula \eqref{eq:23.3.8} for general winding number. Here we will
simply calculate the coefficient by studying the 't~Hooft--Polyakov
monopole~\hyperref[ch23-ref-2]{[2]}, in which the fields have unit winding
number.

As we saw in the previous section, the `identity' (as opposed to the unit)
element $\mathcal I$ of $\pi_2(SU(2)/U(1))$, which corresponds to the
element `one' of $\mathbb Z$, consists of configurations in which the
two-sphere $\mathcal S$ at infinity is mapped once (with positive
Jacobian) into the sphere described by $\hat\phi_n$. As a representative
of this homotopy class, we may take a configuration in which at infinity
$\phi_n$ is in the same direction as $\mathbf x$. To construct this
configuration let us impose a symmetry under joint rotations on
$\boldsymbol\phi=\{\phi_1,\phi_2,\phi_3\}$ and $\mathbf x$, as well as
parity conservation, and make the ansatz:
\begin{equation}
 \phi_n=\hat x_n\langle\phi\rangle F(r),
 \tag{23.3.9}\label{eq:23.3.9}
\end{equation}
\begin{equation}
 A_{ni}=\frac{\epsilon_{nil}\hat x_l}{er}G(r).
 \tag{23.3.10}\label{eq:23.3.10}
\end{equation}
There is an important similarity between this field configuration and
that for a vortex line in a superconductor. The solution
$\phi=\pm\ell\varphi/2e$ for the Goldstone boson field in a vortex line
found in \hyperref[sec:21.6]{Section 21.6} shows that, although gauge
invariance and rotational invariance are spontaneously broken, the vortex
line solution is invariant under the combination of a global gauge
transformation for which $\phi\to\phi+\Lambda$ and a rigid rotation
$\varphi\to\varphi\pm2e\Lambda/\ell$. Similarly, a monopole solution of
the form \eqref{eq:23.3.9}--\eqref{eq:23.3.10} is not invariant under
rotations or gauge transformations, but it is invariant under the
combination of a rigid three-dimensional spatial rotation and an equal
global $SO(3)$ gauge transformation.

As already mentioned, in order for the integral of $V(\phi_n\phi_n)$ to
converge it is necessary for $\phi_n\phi_n$ to approach
$\langle\phi\rangle^2$ as $r\to\infty$, so in this limit $F(r)\to1$. To
derive the limiting behavior of $G(r)$, note that the covariant derivative
of the scalar field is
\begin{equation}
 D_i\phi_n=\pm\langle\phi\rangle\left[
 (1-G(r))(\delta_{ni}-\hat x_n\hat x_i)\frac{F(r)}r
 +\hat x_n\hat x_iF'(r)\right],
 \tag{23.3.11}\label{eq:23.3.11}
\end{equation}
so the scalar term in the Hamiltonian density is
\begin{equation}
 \frac12(D_i\phi_n)^2
 =\langle\phi\rangle^2\left[
 \frac{F^2(1-G)^2}{r^2}+\frac{F'^2}{2}\right].
 \tag{23.3.12}\label{eq:23.3.12}
\end{equation}
For this to have a finite integral it is necessary also that $G(r)\to1$
and $F'(r)\to0$ as $r\to\infty$. Finally, the field strength
\eqref{eq:23.3.2} is
\begin{equation}
 F_{nij}=\frac{\epsilon_{ijk}}e\left[
 -\frac1rG'(r)(\delta_{nk}-\hat x_n\hat x_k)
 -\frac1{r^2}(2G(r)-G^2(r))\hat x_n\hat x_k\right],
 \tag{23.3.13}\label{eq:23.3.13}
\end{equation}
so the Yang--Mills term in the Hamiltonian density is
\begin{equation}
 \frac14(F_{nij})^2=\frac1{e^2}\left[
 \frac{G'^2}{r^2}+\frac{(2G-G^2)^2}{2r^4}\right].
 \tag{23.3.14}\label{eq:23.3.14}
\end{equation}
This has an integral that converges at large distances as long as $G'(r)$
vanishes sufficiently rapidly as $r\to\infty$.

We can use these results to calculate the magnetic monopole moment of this
configuration. Eq.~\eqref{eq:23.3.13}, together with the limits
$G(\infty)=1$ and $G'(\infty)=0$, shows that for $r\to\infty$, we have
\begin{equation}
 F_{nij}\longrightarrow
 -\frac{\epsilon_{ijk}\hat x_k\hat x_n}{er^2}.
 \tag{23.3.15}\label{eq:23.3.15}
\end{equation}
Since $D_i\phi_n$ vanishes rapidly as $r\to\infty$, the magnetic part of
the field strength tensor in this gauge is given for $r\to\infty$ by the
first term in Eq.~\eqref{eq:23.3.5}, so at large distances the magnetic
field becomes
\begin{equation}
 B_i\equiv\frac12\epsilon_{ijk}\mathcal F_{jk}
 \longrightarrow\frac12\epsilon_{ijk}\hat\phi_nF_{njk}
 \longrightarrow\frac{\hat x_i}{er^2}.
 \tag{23.3.16}\label{eq:23.3.16}
\end{equation}
Thus \emph{this configuration has magnetic monopole moment $g=1/e$}.
According to the general argument above, the magnetic moment of a
configuration with winding number $\nu$, corresponding to the element
$\nu$ of $\mathbb Z$, is then
\begin{equation}
 g_\nu=\nu/e.
 \tag{23.3.17}\label{eq:23.3.17}
\end{equation}

It requires a detailed numerical calculation to find the stable
configuration that minimizes the integral of the Hamiltonian density
\eqref{eq:23.3.4}. However there is a limiting case where an analytic
solution is available. To see this, it is useful first to derive a general
lower bound due to Bogomol'nyi~\hyperref[ch23-ref-10]{[10]} on the
monopole energy for a given magnetic monopole moment $g$. Note that
\eqref{eq:23.3.4} may be written
\begin{equation}
 \mathcal H=\frac14(F_{nij}\mp\epsilon_{ijk}D_k\phi_n)^2
 \mathbin{\pm}\frac12\epsilon_{ijk}F_{nij}D_k\phi_n
 +V(\phi_n\phi_n).
 \tag{23.3.18}\label{eq:23.3.18}
\end{equation}
Using the Bianchi identity \eqref{eq:15.3.9}, the second term may be
written
\[
 \mathbin{\pm}\frac12\epsilon_{ijk}F_{nij}D_k\phi_n
 =\mathbin{\pm}\frac12\epsilon_{ijk}D_k(F_{nij}\phi_n)
 =\mathbin{\pm}\frac12\epsilon_{ijk}\partial_k(F_{nij}\phi_n),
\]
so its integral is given by the magnetic monopole moment $g$
\[
 \mathbin{\pm}\frac12\epsilon_{ijk}\int d^3x\,
 F_{nij}D_k\phi_n
 =\mathbin{\pm}\langle\phi\rangle\int\mathbf B\cdot d\mathbf A
 =\mathbin{\pm}4\pi\langle\phi\rangle g.
\]
Since every other term in $\mathcal H$ is positive, we have a general
lower bound on the energy of a configuration with magnetic monopole
moment $g$
\begin{equation}
 E=\int d^3x\,\mathcal H
 \geq4\pi\langle\phi\rangle|g|.
 \tag{23.3.19}\label{eq:23.3.19}
\end{equation}
For $g=\pm1/e$, this gives an energy
$E\geq4\pi\langle\phi\rangle/e$, which for small coupling constant $e$ is
much greater than the corrections due to quantum fluctuations, which are
at most of order $\langle\phi\rangle$. This is why we can take such a
classical configuration seriously as the leading term in a perturbation
expansion.

Now, it is tempting to try to minimize the energy for a given magnetic
monopole moment by setting the first term in Eq.~\eqref{eq:23.3.18} equal
to zero, so that
\begin{equation}
 F_{nij}=\pm\epsilon_{ijk}D_k\phi_n,
 \tag{23.3.20}\label{eq:23.3.20}
\end{equation}
but in general this does not lead to a configuration at which the energy
is stationary. The condition that the energy should be stationary with
respect to variations in the scalar field is the field equation
\begin{equation}
 D_kD_k\phi_n=2\phi_nV'(\phi_n\phi_n),
 \tag{23.3.21}\label{eq:23.3.21}
\end{equation}
while Eq.~\eqref{eq:23.3.20} together with the Bianchi identity
\eqref{eq:15.3.9} would imply that $D_kD_k\phi_n=0$. This argument
suggests that in the special case where $V(\phi_n\phi_n)$ is very small,
it will be possible nearly to reach the lower bound \eqref{eq:23.3.19},
and hence minimize the energy for a given magnetic monopole moment, by
imposing the condition \eqref{eq:23.3.20}. (Where $V$ is the quartic
polynomial $\lambda(\phi_n\phi_n-\langle\phi\rangle^2)^2$, the assumption
that $V$ is small means that $\lambda\ll e^2$, as in a type I
superconductor.) Such stable configurations were studied in this way by
Bogomol'nyi~\hyperref[ch23-ref-10]{[10]}. They had been found earlier by
Prasad and Sommerfield~\hyperref[ch23-ref-14]{[14]} without direct use of
Eq.~\eqref{eq:23.3.20}, and are usually called \emph{BPS monopoles}.

The Bogomol'nyi condition \eqref{eq:23.3.20} provides first-order
differential equations for $F(r)$ and $G(r)$, which are much easier to
solve than the second-order field equations derived directly from the
condition that the energy should be stationary. Using
Eqs.~\eqref{eq:23.3.11} and \eqref{eq:23.3.13}, the terms in
Eq.~\eqref{eq:23.3.20} proportional to
$\epsilon_{ijk}[\delta_{kn}-\hat x_k\hat x_n]$ and
$\epsilon_{ijk}\hat x_k\hat x_n$ respectively yield the differential
equations
\begin{equation}
 e\langle\phi\rangle F(1-G)=G',
 \tag{23.3.22}\label{eq:23.3.22}
\end{equation}
\begin{equation}
 e\langle\phi\rangle r^2F'=G(2-G).
 \tag{23.3.23}\label{eq:23.3.23}
\end{equation}
With the boundary condition that $F(r)\to1$ and $G(r)\to1$ as
$r\to\infty$, these equations have the solution
\begin{equation}
 F=\coth\rho-\frac1\rho,\qquad
 G=1-\frac{\rho}{\sinh\rho},
 \tag{23.3.24}\label{eq:23.3.24}
\end{equation}
where $\rho=e\langle\phi\rangle r$. Note in particular that the field
$\phi_n$ given by Eqs.~\eqref{eq:23.3.9} and \eqref{eq:23.3.24} vanishes
for $r\to0$, so as remarked in \hyperref[sec:23.1]{Section 23.1}, the
$SU(2)$ symmetry is restored at the center of the monopole.

Let's now return to the case of a potential $V$ of arbitrary strength. The
't~Hooft--Polyakov monopole is stable, because are no configurations of
smaller topological quantum number into which it could decay.
Configurations of higher magnetic monopole moment are generally
unstable~\hyperref[ch23-ref-10]{[10]}. There are also interesting
configurations with both magnetic monopole moments and electric charge,
known as \emph{dyons}~\hyperref[ch23-ref-15]{[15]}.

There is another way of understanding the value $1/e$ of the
't~Hooft--Polyakov magnetic monopole moment, which goes back to the
original work of Dirac on magnetic monopoles \hyperref[ch23-ref-16]{[16]}.
As mentioned earlier, instead of the gauge we have been using, we can make
a gauge transformation
$\phi_n\to R_{nm}(\hat{\mathbf x})\phi_m$ that rotates $\phi_n$ to point
in a fixed direction $\hat{\mathbf v}$, for instance the three-direction.
Then the field strength
$B_{nk}\equiv\frac12\epsilon_{ijk}F_{nij}$ is transformed into
$R_{nm}B_{mk}$, which approaches
$\hat v_n\hat x_k/er^2$ for $r\to\infty$, so here we do not have to
project this on a local unbroken symmetry direction. The price we need to
pay for these conveniences is that the gauge transformation is singular;
the rotation that takes a vector in the $\hat{\mathbf x}$ direction into
some fixed direction $\hat{\mathbf v}$ is
\begin{align}
 R(\hat{\mathbf x};\hat{\mathbf v})={}&
 \mathbf 1-(1-\hat{\mathbf v}\cdot\hat{\mathbf x})
 \hat{\mathbf v}\hat{\mathbf v}^{T}
 +\hat{\mathbf v}
 \bigl(\hat{\mathbf x}-(\hat{\mathbf x}\cdot\hat{\mathbf v})
 \hat{\mathbf v}\bigr)^{T}
 \notag\\
 &+\bigl(\hat{\mathbf x}-(\hat{\mathbf x}\cdot\hat{\mathbf v})
 \hat{\mathbf v}\bigr)\hat{\mathbf v}^{T}
 \notag\\
 &-\frac{
 \bigl(\hat{\mathbf x}-(\hat{\mathbf x}\cdot\hat{\mathbf v})
 \hat{\mathbf v}\bigr)
 \bigl(\hat{\mathbf x}-(\hat{\mathbf x}\cdot\hat{\mathbf v})
 \hat{\mathbf v}\bigr)^{T}}
 {1+\hat{\mathbf x}\cdot\hat{\mathbf v}},
 \tag{23.3.25}\label{eq:23.3.25}
\end{align}
which is singular at $\hat{\mathbf x}=-\hat{\mathbf v}$. This $R$ is not
unique; for instance, we could perform the rotation
$R(\hat{\mathbf x};-\hat{\mathbf v})$ that takes $\hat{\mathbf x}$ into
the $-\hat{\mathbf v}$ direction, followed by a fixed rotation of
$180^\circ$ around some axis perpendicular to $\hat{\mathbf v}$, but this
would become singular at $\hat{\mathbf x}=+\hat{\mathbf v}$. To avoid the
singularities, we have to adopt different gauges in different regions; for
instance, with $\hat{\mathbf v}$ in the three-direction, we can use a gauge
for $0<\theta<\theta_0$ that is singular at $\theta=\pi$ and a gauge for
$\theta_0<\theta<\pi$ that is singular at $\theta=0$, where $\theta_0$ is
an arbitrary angle with $0<\theta_0<\pi$, often taken as $\pi/2$.
Everywhere except at $\theta=0$ and $\theta=\pi$ the magnetic field will
be given at large distances by
$\mathbf B\to g\hat{\mathbf x}/r^2$, where $g$ is the magnetic monopole
strength. This can be written as the curl of a vector potential whose only
non-vanishing component is in the azimuthal $\varphi$ direction. For
$0<\theta<\theta_0$ we must take
$A_\varphi=g(1-\cos\theta)/r\sin\theta$, which is only singular at
$\theta=\pi$, while for $\theta_0<\theta<\pi$ we must take
$A_\varphi=-g(1+\cos\theta)/r\sin\theta$, which is only singular at
$\theta=0$. The difference $\Delta\mathbf A$ between these two vector
potentials is a gradient $\boldsymbol\nabla\Lambda$ with
$\Lambda=2g\varphi$, which of course does not affect the magnetic field
for $0<\theta<\pi$, but could affect the dynamics of charged fields. A
gauge transformation with $\Lambda=2g\varphi$ will change a field of
charge $q$ by a factor $\exp(2iqg\varphi)$, which is not single-valued
unless $2qg$ is an integer. This is the Dirac quantization condition; the
existence of any magnetic monopole with monopole moment $g$ would require
all electric charges to be integer multiples of $(2g)^{-1}$. For the
't~Hooft--Polyakov monopole this condition is automatically satisfied,
because here $g=1/e$, and all charges in the Georgi--Glashow model are
integer multiples of $e/2$.

The Georgi--Glashow model was ruled out as a theory of weak and
electromagnetic interactions by the discovery of neutral currents, but
magnetic monopoles are expected to occur in other theories, where a simply
connected group $G$ is spontaneously broken not to $U(1)$, but to some
subgroup $H'\cdot U(1)$, where $H'$ is simply connected. (According to
\hyperref[app:23.B]{Appendix B of this chapter}, for simply connected
groups $G$ we have $\pi_2(G/H)=\pi_1(H)$, which for
$H=H'\cdot U(1)$ equals
$\pi_1(H')\cdot\pi_1(U(1))=\pi_1(U(1))=\mathbb Z$.) There are no
monopoles produced in the spontaneous breaking of the gauge group
$SU(2)\cdot U(1)$ of the standard electroweak theory, which is not simply
connected. (About this, more below.) But we do find monopoles when the
simply connected gauge group $G$ of theories of unified strong and
electroweak interactions, such as $SU(4)\cdot SU(4)$ or $SU(5)$ or
$\operatorname{Spin}(10)$, is spontaneously broken to the gauge group
$SU(3)\cdot SU(2)\cdot U(1)$ of the standard model. (See
\hyperref[sec:21.5]{Section 21.5}.) The monopoles in this case are expected
to have a mass larger by an inverse square gauge coupling constant than
the vector boson masses $M\approx10^{15}$--$10^{16}\ \mathrm{GeV}$
produced by this symmetry breaking. Such monopoles would have been
produced when the universe underwent a phase transition in which $G$ was
spontaneously broken to $SU(3)\cdot SU(2)\cdot U(1)$, at a temperature
$T$ of order $M$.

This poses a problem for some cosmological
models~\hyperref[ch23-ref-17]{[17]}. The scalar fields before this phase
transition would have necessarily been uncorrelated at distances larger
than the horizon distance, the furthest distance that light could have
travelled since the initial singularity. At an early time $t$ in standard
cosmological theories~\hyperref[ch23-ref-18]{[18]} the horizon distance is
of order
$t\approx(G_NT^4)^{-1/2}$ (where
$G_N\approx(10^{19}\ \mathrm{GeV})^{-2}$ is Newton's constant), so the
number density of monopoles produced at this time would have been of
order $t^{-3}\approx(G_NM^4)^{3/2}$, which is smaller than the photon
density $M^3$ at $T\approx M$ by a factor of order $(G_NM^2)^{3/2}$. For
$M\approx10^{15}\ \mathrm{GeV}$ this factor is of order $10^{-6}$. If
monopoles did not find each other to annihilate, then this ratio would
remain roughly constant to the present, but with at least $10^9$ microwave
background photons per nucleon today, this would give at least $10^3$
monopoles per nucleon, in gross disagreement with what is observed. This
potential paradox was one of the factors leading to inflationary
cosmological models~\hyperref[ch23-ref-19]{[19]}, in which there was a
period of exponential expansion, which if it occurred before the
monopoles were produced would have greatly extended the horizon, and if
it occurred after the production of monopoles (but before a period of
reheating) would have greatly diluted the monopole density.

The discovery of monopoles of any sort would create opportunities for the
observation of remarkable phenomena, including the existence of
fermion--monopole configurations of fractional fermion
number~\hyperref[ch23-ref-20]{[20]}, and the violation of baryon
conservation in fermion--monopole
scattering~\hyperref[ch23-ref-21]{[21]}.

\begin{center}
$* \quad * \quad *$
\end{center}

In the above discussion we have considered only monopoles associated with
the spontaneous breakdown of a simply-connected gauge group $G$. This
raises a question. For every Lie group $G$, whether simply connected or
not, there is a simply-connected group $\overline G$ with the same Lie
algebra, known as its \emph{covering group}. (For examples, see
\hyperref[sec:2.7]{Section 2.7}.) Any non-simply connected group has fewer
representations than its covering group (for instance, the
doubly-connected groups $SO(n)$ have only scalar, vector, and tensor
representations, while their covering groups $\operatorname{Spin}(n)$ also
have spinor representations). If a theory does not happen to involve
fields that belong to the extra representations of the covering group,
then are we are free to consider the gauge group of the theory to be
either the non-simply connected group $G$ or its covering group
$\overline G$? In particular, does the menu of possible monopoles depend
on whether the theory contains only fields transforming as representations
of a non-simply connected group $G$, or additional fields that would only
furnish representations of its covering group $\overline G$? For instance,
in the original Georgi--Glashow model~\hyperref[ch23-ref-12]{[12]} the
only scalar fields belonged to a representation of $SO(3)$, the
three-vector, but there were fermions that belonged to spinor
representations of the covering group
$\operatorname{Spin}(3)=SU(2)$. Would the allowed values of the magnetic
monopole moment change if we added scalar fields belonging to the spinor
representations of $SU(2)$? Would it change if we removed the fermions?

The answer is that the menu of possible monopoles does not depend on
whether we say that the gauge group is a non-simply connected group $G$ or
its covering group $\overline G$, and therefore is unaffected when we add
or remove fields belonging to representations of $\overline G$ that are
not representations of $G$. As we saw in
\hyperref[sec:23.1]{Section 23.1}, in general the topologically stable
monopole-like configurations are classified according to the elements of
$\pi_2(G/H)$. According to a result quoted in
\hyperref[app:23.B]{Appendix B of this chapter}, this homotopy group
consists of those elements of $\pi_1(H)$ that correspond to the trivial
element of $\pi_1(G)$ when $H$ is embedded in $G$. But if we replace $G$
with its covering group $\overline G$, we also replace $H$ with a
different subgroup $H'$, because some of the loops in $H$ do not return to
the base point when $H$ is embedded in $\overline G$. These are just the
loops that do not become trivial when $H$ is embedded in $G$, so
$\pi_2(G/H)=\pi_1(H')$, and thus as far as monopoles are concerned we
could just as well say that the gauge group is $\overline G$ rather than
$G$.

For instance, as far as the scalar fields are concerned, the gauge group
of the Georgi--Glashow model might be considered to be the
doubly-connected $SO(3)$ rather than its simply-connected covering group
$SU(2)$. The unbroken subgroup is then $SO(2)$, in which we identify
transformations that differ by a $360^\circ$ rotation. Thus
$\pi_1(SO(2))$ includes loops that extend from the unit element to a
$360^\circ$ rotation, which would not be loops when $SO(3)$ is embedded
in $SU(2)$. But $\pi_2(SO(3)/SO(2))$ is not the same as
$\pi_1(SO(2))$, but rather excludes the loops that are homotopically
non-trivial when $SO(3)$ is embedded in $SU(2)$, which are just the loops
that extend from the unit element to a $360^\circ$ rotation, so
$\pi_2(SO(3)/SO(2))$ is the $U(1)$ subgroup of $SU(2)$, just as if the
gauge group were taken to be $SU(2)$ from the beginning.

It is convenient always to consider the gauge group $G$ associated with
any semi-simple gauge algebra to be the simply-connected covering group,
so that we can use the simple result that $\pi_2(G/H)=\pi_1(H)$. As we
have just seen, the connectivity properties of $H$ are then fixed by its
embedding in $G$, or more precisely, by the embedding of the Lie algebra
of $H$ in the Lie algebra of $G$. For instance, an $SU(3)$ gauge algebra
might be spontaneously broken either into an $SU(2)$ subalgebra, under
which the defining representation of $SU(3)$ transforms as a doublet plus
a singlet, or into an $SO(3)$ subalgebra, under which the defining
representation of $SU(3)$ transforms as a three-vector. In the first case
we do not have the option of considering the unbroken subgroup to be
$SO(3)$; since $\pi_1(SU(2))=0$, there are no monopoles here. In the
second case the unbroken gauge group must be regarded as $SO(3)$, not
$SU(2)$, so the theory does have monopole-like configurations, classified
according to the elements of $\pi_1(SO(3))=\mathbb Z_2$. The nature of
the unbroken subalgebra of $H$ and its embedding in the gauge algebra $G$
may be dynamically affected by the variety of field types that we
introduce in the Lagrangian, but once the algebra of $H$ and its embedding
in the algebra of $G$ are fixed, the menu of monopoles is otherwise
entirely unaffected by the variety of fields in the theory.

In particular, the argument that led to the Dirac quantization condition
shows that, in any theory in which a Lie algebra $G$ is spontaneously
broken to a subalgebra including the electric charge operator, the allowed
magnetic monopole moments are integer multiples of the reciprocal of the
smallest electric charge that appears in the representations of the
covering group of $G$, whether or not there is actually any particle that
carries that charge in the theory. If the algebra of $G$ itself contains a
$U(1)$ generator, then we must consider the covering group of this
$U(1)$, which is the non-compact group of translations along the real
line. If this $U(1)$ generator appears as a term in the electric charge
operator, as in the standard electroweak theory, then there is no minimum
electric charge in the representations of the covering group, and hence
no monopoles.
